% =====================================================================
%  Introduction — Multi-Period OPF with Parallel Temporal Decomposition
%  Drop-in section for an IEEE Transactions paper.
%  Requires: amsmath, amssymb, cite (or natbib)
% =====================================================================

\section{Introduction}
\label{sec:intro}

The rapid proliferation of distributed photovoltaic (PV) generation and battery
energy storage systems (BESS) in distribution networks has fundamentally changed
the nature of the optimal power flow (OPF) problem faced by distribution system
operators. When batteries are co-located with PV, each battery's state of charge
(SOC) at the end of one hour determines the feasible operating range at the next,
coupling consecutive timesteps and converting what was a sequence of independent
single-period problems into a single temporally coupled multi-period program.
Day-ahead scheduling of a feeder that hosts hundreds of DER–BESS pairs over a
24-hour horizon at hourly resolution requires solving this coupled program in its
entirety: truncating the horizon, ignoring future prices, or approximating the
coupling with a short look-ahead window sacrifices the energy-arbitrage opportunity
that the battery was installed to exploit. The challenge is therefore to solve
the full-horizon, temporally coupled multi-period OPF accurately and fast enough
to be operationally useful on realistic, large-scale unbalanced three-phase feeders.

\subsection*{Background and related work}

\paragraph{Convex relaxations for three-phase distribution OPF.}
The branch-flow model (BFM) of Baran and Wu~\cite{BaWu89} casts the power-flow
equations in terms of branch currents and complex voltages, exposing a
second-order cone structure that Jabr~\cite{Jabr06} and Farivar and
Low~\cite{FaLo13} exploit to obtain a tractable convex relaxation.
For balanced radial networks the SOCP relaxation is exact under mild
conditions~\cite{FaLo13}; its linearization, LinDistFlow~\cite{BaWu89}, trades
loss accuracy for a simple linear program. These single-phase results have been
extended to the full unbalanced three-phase BFM~\cite{GaLo14,ZhLo16}, giving
operators a hierarchy of convex surrogates: a lossless LP at the bottom, an
SOCP relaxation in the middle, and the non-convex AC model at the top.
None of these models, however, is used in isolation for multi-period scheduling:
the SOCP relaxation may be inexact for unbalanced networks with mutual coupling,
and the non-convex AC model is intractable for large feeders over a full day's
horizon.

\paragraph{Iterative SOCP for AC-accurate solutions.}
A line of work addresses the relaxation gap by iteratively tightening the conic
constraints until an AC-consistent operating point is
recovered~\cite{NiGr16,BoBu19}. The key idea is to add directional cuts that
penalize or exclude the infeasible portion of the cone after each solve, so the
sequence of SOCP iterates converges to the boundary of the feasibility set.
Our prior work formulated such an iterative SOCP (MPISOCP) for the multi-period
three-phase BFM, establishing that the scheme converges in a small number of
iterations (8--9 on the IEEE 123-bus feeder) and that its objective matches the
full nonlinear AC solution to within \$0.14 over a 24-hour
horizon~\cite{OurISCOPRef}. The MPISOCP is the OPF engine used inside each
temporal window of the decomposition proposed in this paper.

\paragraph{Multi-period OPF with BESS in distribution networks.}
Incorporating BESS into multi-period OPF introduces the SOC recursion as the
sole inter-temporal coupling, but that coupling is enough to make the monolithic
program grow super-linearly with the horizon length. Nazir and
Racherla~\cite{NaRa20} solve the multi-period three-phase OPF on unbalanced
feeders using an SOCP relaxation, dispatching batteries to minimize losses
over a day. Crucially, they adopt a \emph{receding-horizon} (model-predictive
control) strategy: at each hour a short look-ahead window is optimized, only
the first step is executed, and the window shifts forward. This is operationally
simple but inherently myopic — the battery cannot plan more than a few hours
ahead, so it forfeits the global energy-arbitrage opportunity that spans the
full midday PV surplus and the late-afternoon price peak. Avramidis and
Capitanescu~\cite{AvCa20} develop a comprehensive multi-period OPF framework for
low-voltage distribution networks with storage, but their centralized formulation
does not scale beyond relatively small feeders.
Blasi \textit{et al.}~\cite{BlFe22} and similar works solve multi-period active
distribution network OPF centrally for moderate-size test systems, without
addressing the scalability barrier that arises on large, heavily unbalanced
feeders.

\paragraph{Temporal decomposition for multi-period power system optimization.}
The idea of breaking a long-horizon optimization into shorter, coupled sub-problems
traces back to Benders decomposition~\cite{Be62} and its multi-stage
generalization, stochastic dual dynamic programming (SDDP)~\cite{PePi91}.
In SDDP, each stage is a two-stage problem: it receives a forward state (here the
SOC) from its predecessor and returns a piecewise-linear cost-to-go cut to its
successor. The cuts are valid globally because each stage's value function is
convex in its initial state. This sequential backward-then-forward sweep has been
applied to battery scheduling in microgrids~\cite{BhKh16}, and Lan
\textit{et al.}~\cite{LaZh22} develop a fast SDDP for economic dispatch in
distribution systems with storage. However, these works adopt simplified
linearized or single-phase network models and do not address the three-phase
unbalanced AC physics that govern realistic distribution feeders. Giuntoli
\textit{et al.}~\cite{GiSu21} apply a nested Benders scheme to multi-period OPF
for distribution grids with BESS, but their formulation is sequential — each
stage must wait for all upstream cuts — and is tested only on small systems.
Agarwal and Pileggi~\cite{AgPi21} use differential dynamic programming (DDP) to
decompose large-scale multi-period OPF with energy storage, demonstrating
scalability at bulk-system scale; their formulation targets balanced, largely
transmission-scale networks and does not incorporate unbalanced three-phase
distribution physics.

Beyond temporal decomposition, spatial ADMM-based methods have been widely
studied for multi-area distribution OPF~\cite{PeLo16,ErPa14}. These decompose
across network areas (spatial decomposition) rather than across time, and
they do not address the temporal coupling introduced by BESS. A combined
spatial–temporal ADMM approach~\cite{PiBe20} has been explored on a small
33-bus low-voltage feeder but does not recover global optimality guarantees and
does not scale to nearly 10,000-node systems.

\paragraph{The gap.}
The literature reveals three orthogonal capabilities that have never been
combined in a single framework:

\begin{enumerate}
  \item \textbf{AC-accurate three-phase unbalanced physics.} Receding-horizon
        and centralized SOCP works~\cite{NaRa20,AvCa20} handle the full network
        model but sacrifice global optimality or scalability.
  \item \textbf{Provably optimal temporal decomposition.} SDDP-style and nested
        Benders works~\cite{LaZh22,GiSu21,AgPi21} provide optimality guarantees
        but restrict the network model to single-phase, linearized, or balanced
        formulations.
  \item \textbf{Parallel dispatch of all temporal windows.} Classical SDDP and
        nested Benders sweep the window chain sequentially; the potential to
        solve all windows simultaneously per iteration — and thereby reduce
        wall-clock time proportionally to the available cores — is unexploited
        in the distribution OPF context.
\end{enumerate}

\noindent The barriers to combining them are non-trivial. The nonconvex
three-phase BFM makes each window subproblem hard; the ISOCP tightening that
recovers AC accuracy requires multiple inner iterations and, if repeated at every
outer Benders iteration, would dominate the total runtime. Parallel dispatch
requires all windows to be self-sufficient at each outer iteration, using only
information available from the previous iteration (forward SOC states and
backward dual cuts), so no window may block on a neighbor within an iteration.
The present paper resolves all three issues simultaneously.

\subsection*{Contributions}

This paper presents \emph{DDDP--OTD}, a Parallel Dual Dynamic Decomposition
over Time for multi-period OPF on unbalanced three-phase distribution networks
with DERs and BESS. The specific contributions are as follows.

\begin{enumerate}

\item \textbf{Parallel multi-cut Benders temporal decomposition (DDDP--OTD).}
The 24-hour scheduling horizon is partitioned into $P$ disjoint contiguous
windows. Within each outer iteration, all $P$ window subproblems are solved
simultaneously: each consumes the terminal SOC forwarded by its predecessor and
the optimality cut emitted by its successor at the \emph{previous} iteration,
so no in-iteration synchronization is required. A monotone-nondecreasing global
lower bound and an \emph{exact} primal upper bound are available at every
iteration, and the gap between them closes at a rate independent of $P$.

\item \textbf{Deferred AC tightening.}
Each window subproblem is solved by the MPISOCP scheme, but the full
iterative cone tightening is \emph{not} rerun at every outer Benders iteration.
Instead, each iteration admits only a fixed number of cheap inner SOCP cuts,
and the residual conic gap is closed by a single complete ISOCP sweep once
DDDP--OTD has converged. This amortizes the costly AC recovery over the entire
horizon rather than paying it per window per Benders iteration, so a near-exact
AC-OPF schedule is obtained for essentially the cost of a single pass through the
relaxed decomposition.

\item \textbf{Demonstrated scalability on large unbalanced feeders.}
On the IEEE 123-bus feeder, DDDP--OTD reduces the centralized ISOCP solve time
from $23.8$~s to $11.4$~s ($2.1\times$) at $P\!=\!8$ with an optimality gap
below $0.002\%$ and an objective that matches the nonlinear AC solution to
\$0.14. On a 2752-bus reduction of the nearly 10,000-node IEEE 9500-node feeder,
the centralized solve of $977.6$~s ($16.3$~min) is reduced to $272.3$~s
($3.6\times$) at $P\!=\!6$, with a gap below $0.006\%$ --- the first parallel
temporal decomposition of an AC-accurate multi-period OPF demonstrated at this
scale on an unbalanced feeder.

\item \textbf{End-to-end validation against unbalanced power flow.}
Every recovered schedule is replayed hour-by-hour through a full unbalanced
OpenDSS power flow. The ISOCP-within-DDDP schedule matches OpenDSS to within
$3.1$~kW and $16.5$~kVAr at the substation and $1.3\times10^{-4}$~p.u.\ in
voltage on the 9500-node feeder, confirming that the schedule is physically
deliverable at scale.

\end{enumerate}

\subsection*{Paper organization}

The remainder of the paper is organized as follows.
Section~\ref{sec:copf} formulates the centralized multi-period three-phase OPF,
presents the nonlinear and linear models, and develops the MPISOCP convexification
scheme. Section~\ref{sec:dddp} introduces the DDDP--OTD decomposition: window
construction, the forward–backward exchange of SOC states and dual cuts, global
bound tracking, and the parallel algorithm.
Section~\ref{sec:results} reports results on the IEEE 123-bus and 9500-node
feeders, including the $\gamma$-parameter study, OpenDSS validation, and the
partition-count trade-off.
Section~\ref{sec:conclusion} concludes with directions for future work.
